\documentclass[12pt]{article}
\usepackage[utf8]{inputenc}
\usepackage[T1]{fontenc}
\usepackage[margin=1in]{geometry}
\usepackage{setspace}
\usepackage{xcolor}
\usepackage[colorlinks=true, urlcolor=blue!60!black, linkcolor=blue!60!black]{hyperref}
\setlength{\parindent}{2em}
\setlength{\parskip}{1.25em}
\renewcommand{\baselinestretch}{1.0}
\title{\textit{Hayburn's Case}, 2 U.S. (2 Dall.) 409 (1792)}
\author{Case Brief}
\date{}
\begin{document}
\maketitle

\section{Case Metadata}

\begin{itemize}
  \item \textbf{Case Name:} \textit{Hayburn's Case}
  \item \textbf{Citation:} \href{https://supreme.justia.com/cases/federal/us/2/409/}{2 U.S. (2 Dall.) 409 (1792)}
  \item \textbf{Court and Dates:} Supreme Court of the United States; argued August 1792 Term (August 11, 1792); held under advisement and mooted by Act of February 28, 1793
  \item \textbf{Chapter:} Chapter 1: Judicial Review and the Federal Courts
  \item \textbf{Vote:} 3--3 split denying Attorney General Randolph's \textit{ex officio} mandamus motion; no final merits vote after Congress amended the statute (five of the six sitting Justices---Jay, Cushing, Wilson, Blair, and Iredell---expressed their constitutional views in circuit-court letters appended by reporter Alexander Dallas at 2 U.S. at 410--414 n.*)
  \item \textbf{Author:} No formal merits opinion (procedural report by Alexander Dallas with extensive footnote reprinting the circuit-court letters to President George Washington)
  \item \textbf{Constitutional and Statutory Provisions:} U.S. Const. art. III, \S~1 (Judicial Vesting Clause); art. III, \S~2, cl. 1 (Case-or-Controversy Requirement); Invalid Pensions Act of March 23, 1792, ch. 11, 1 Stat. 243
  \item \textbf{Primary Source Links:} \href{https://supreme.justia.com/cases/federal/us/2/409/}{Justia U.S. Supreme Court Center}; \href{https://www.law.cornell.edu/constitution-conan/article-3/section-2/clause-1/advisory-opinion-doctrine-and-practice}{Cornell LII Constitution Annotated}; \href{https://www.fjc.gov/history/cases/cases-that-shaped-the-federal-courts/hayburns-case}{Federal Judicial Center}
\end{itemize}

\section{Facts and Procedural History}

On March 23, 1792, Congress enacted the Invalid Pensions Act to provide disability pensions for wounded Revolutionary War veterans. The statute directed disabled veterans to apply to the federal circuit courts, which were instructed to examine medical proof of disability, determine each applicant's eligibility and pension rate, and certify their findings to the Secretary of War. Crucially, however, the statute made the courts' determinations non-final: the Secretary of War could withhold a veteran's name from the pension list whenever he had ``cause to suspect imposition or mistake,'' and Congress retained ultimate authority to revise or reject the pension awards.

Because Supreme Court Justices then ``rode circuit'' alongside district judges, five of the six sitting Justices confronted the Act immediately on the circuit benches and objected that it violated Article III separation of powers. In April 1792, the Circuit Court for the District of Pennsylvania (Justices James Wilson and John Blair, and District Judge Richard Peters) refused to hear the pension petition of Revolutionary War veteran William Hayburn and sent a joint letter to President George Washington explaining that the Act was unconstitutional. Circuit courts in New York (Chief Justice John Jay, Justice William Cushing, and Judge James Duane) and North Carolina (Justice James Iredell and Judge John Sitgreaves) likewise wrote to President Washington declaring the judicial assignment unconstitutional, though the New York judges voluntarily agreed to hear claims out of court as non-judicial ``commissioners.''

In August 1792, Attorney General Edmund Randolph moved in the Supreme Court for a writ of mandamus compelling the Pennsylvania Circuit Court to proceed on Hayburn's petition. Randolph initially moved \textit{ex officio} on behalf of the United States without presidential instruction or a private client; the six-Justice Court divided 3--3 on his authority to do so and denied the \textit{ex officio} motion. Randolph immediately renewed the motion as private counsel for William Hayburn. Expressing reluctance to rush a constitutional clash with Congress, the Supreme Court held the motion under advisement until its next term. Before the Court could rule, Congress enacted the Act of February 28, 1793, creating a non-judicial administrative procedure for pension claims and rendering Hayburn's mandamus petition moot.

\section{Issue(s) Presented}

\begin{enumerate}
  \item \textbf{Non-Judicial Duties and Judicial Finality:} May Congress assign Article III federal courts to adjudicate statutory claims whose decisions are subject to revision, suspension, or reversal by executive officers (the Secretary of War) or by Congress itself?
  \item \textbf{Attorney General \textit{Ex Officio} Standing:} May the Attorney General seek a writ of mandamus \textit{ex officio} to compel execution of a federal statute without authorization from the President or representation of an injured private party?
\end{enumerate}

\section{Holding}

\textbf{No formal Supreme Court merits judgment (motion mooted by 1793 amendment); foundational Article III principle established via the circuit-court letters.} Although the Supreme Court never issued a final merits opinion in \textit{Hayburn's Case} (having split 3--3 on Randolph's \textit{ex officio} authority and taken his renewed private motion under advisement until Congress repealed the contested procedure), the unanimous constitutional consensus expressed by five Justices in the reported circuit-court letters established two enduring pillars of Article III: (1) Congress may not assign Article III courts non-judicial duties; and (2) Article III forbids federal courts from rendering advisory or non-final judgments subject to executive or legislative revision.

\section{Reasoning}

The doctrinal reasoning of \textit{Hayburn's Case} resides in the three circuit-court protests transmitted to President Washington and preserved in Alexander Dallas's reporter's footnote (2 U.S. at 410--414 n.*):

\textbf{1. Strict Tripartite Separation of Powers.} All three circuit courts began from the constitutional architecture of Articles I, II, and III: the Constitution divides the federal government into three distinct and independent branches, vesting ``the judicial Power of the United States'' exclusively in Article III courts. Neither Congress nor the President may compel Article III courts to perform administrative or legislative duties that are not ``properly judicial'' and performed ``in a judicial manner.''

\textbf{2. Finality as an Essential Attribute of Judicial Power.} The core constitutional defect of the 1792 Act was its revision mechanism. By empowering the Secretary of War---an executive officer---to set aside a circuit court's findings for suspected ``imposition or mistake,'' and by subjecting those findings to congressional override, Congress effectively turned federal courts into subordinate advisors to the political branches. As the Pennsylvania judges emphasized, subjecting judicial decisions to executive and legislative revision was ``radically inconsistent with the independence of that judicial power which is vested in the courts.'' The North Carolina judges went further, explaining that no decision of a federal court could be made liable to reversion or suspension by the Legislature itself, which holds no judicial power apart from impeachment.

\textbf{3. Judicial Independence vs. Individual Service as Commissioners.} While all three circuits agreed that the Act could not bind federal \textit{courts} sitting judicially, they diverged on whether individual judges could voluntarily execute the statute as non-judicial ``commissioners'' designated by official title rather than personal name---a distinction later echoed in separation-of-powers doctrine governing extrajudicial service.

\section{Concurring and Dissenting Opinions}

Because the Supreme Court issued no formal merits opinion (and only an unreported 3--3 division denying Randolph's \textit{ex officio} motion), there were no concurring or dissenting opinions in the modern sense. However, the three circuit-court letters printed in Dallas's footnote reflect notable methodological nuances among the Justices:
\begin{itemize}
  \item \textbf{New York Circuit (Jay, C.J., Cushing, J., and Duane, D.J.):} Held that the statutory duties were non-judicial and could not be imposed on the court, but construed the Act benevolently as designating the judges individually as ``commissioners'' whom they were ``at liberty to accept or decline.'' Out of respect for Congress and Revolutionary veterans, they adjourned court and sat as commissioners between sessions.
  \item \textbf{Pennsylvania Circuit (Wilson, J., Blair, J., and Peters, D.J.):} Adopted the strictest posture, flatly refusing to act on William Hayburn's petition in any capacity because the Act assigned non-judicial business and subjected judicial findings to executive and legislative revision.
  \item \textbf{North Carolina Circuit (Iredell, J., and Sitgreaves, D.J.):} Emphasized that no federal court judgment may be revised or suspended by Congress or the Executive, and expressed serious doubt whether federal judges could constitutionally accept the separate office of ``commissioner'' under the Act without violating the incompatibility of judicial and executive functions.
\end{itemize}

\section{Key Quotations}

\begin{quote}
``That by the Constitution of the United States, the government thereof is divided into three distinct and independent branches, and that it is the duty of each to abstain from, and to oppose, encroachments on either. That neither the Legislative nor the Executive branches can constitutionally assign to the Judicial any duties but such as are properly judicial, and to be performed in a judicial manner.''
--- New York Circuit Court (Jay, C.J., Cushing, J., Duane, D.J.), \href{https://supreme.justia.com/cases/federal/us/2/409/}{2 U.S. (2 Dall.) at 410 n.*}.
\end{quote}

\begin{quote}
``Such revision and control we deemed radically inconsistent with the independence of that judicial power which is vested in the courts; and, consequently, with that important principle which is so strictly observed by the Constitution of the United States.''
--- Pennsylvania Circuit Court (Wilson, J., Blair, J., Peters, D.J.), \href{https://www.law.cornell.edu/constitution-conan/article-3/section-2/clause-1/advisory-opinion-doctrine-and-practice}{2 U.S. (2 Dall.) at 411 n.*}.
\end{quote}

\begin{quote}
``And we beg leave to add, with all due deference, that no decision of any court of the United States can, under any circumstances, in our opinion, agreeable to the Constitution, be liable to a reversion, or even suspension, by the Legislature itself, in whom no judicial power of any kind appears to be vested, but the important one relative to impeachments.''
--- North Carolina Circuit Court (Iredell, J., Sitgreaves, D.J.), \href{https://supreme.justia.com/cases/federal/us/2/409/}{2 U.S. (2 Dall.) at 413 n.*}.
\end{quote}

\begin{quote}
``But the court being divided in opinion on that question, the motion, made ex officio, was not allowed.''
--- Supreme Court Reporter's Entry on Attorney General Randolph's Motion, \href{https://supreme.justia.com/cases/federal/us/2/409/}{2 U.S. (2 Dall.) at 409}.
\end{quote}

\section{Doctrinal Significance and Modern Legacy (Through 2026)}

Despite never producing a binding Supreme Court merits opinion, \textit{Hayburn's Case} is one of the most cited structural precedents in federal courts jurisprudence---predating \textit{Marbury v. Madison} (1803) by eleven years as the first recorded assertion of federal judicial review of an Act of Congress and anchoring three enduring doctrines through 2026:
\begin{itemize}
  \item \textbf{Finality of Article III Judgments:} In \href{https://supreme.justia.com/cases/federal/us/514/211/}{\textit{Plaut v. Spendthrift Farm, Inc.}, 514 U.S. 211 (1995)}, Justice Scalia relied directly on \textit{Hayburn's Case} to invalidate a federal statute requiring Article III courts to reopen final securities-fraud judgments dismissed as time-barred. Article III establishes a ``judicial department'' with the exclusive power to render dispositive judgments that neither Congress nor the Executive may revise retroactively.
  \item \textbf{Prohibition on Advisory Opinions and Non-Judicial Duties:} Alongside the 1793 Correspondence of the Justices, \textit{Hayburn's Case} founded the Article III advisory-opinion bar applied in \href{https://supreme.justia.com/cases/federal/us/219/346/}{\textit{Muskrat v. United States}, 219 U.S. 346 (1911)} (striking down a collusive test-case statute) and reaffirmed in modern standing decisions such as \textit{California v. Texas}, 593 U.S. 659 (2021), which refuse to issue advisory rulings incapable of granting effectual relief. Conversely, in \textit{Mistretta v. United States}, 488 U.S. 361 (1989), the Court invoked the New York circuit judges' ``commissioner'' theory from \textit{Hayburn's Case} to uphold individual federal judges serving voluntarily on the U.S. Sentencing Commission (\href{https://www.law.cornell.edu/constitution-conan/article-3/section-1/inherent-powers-of-federal-courts-issuing-judgments}{Cornell LII Constitution Annotated}).
  \item \textbf{Core Judicial Power and the Administrative State Through 2026:} \textit{Hayburn's Case}'s insistence that Article III judicial power ``cannot be shared with the other branches'' underpins \href{https://supreme.justia.com/cases/federal/us/564/462/}{\textit{Stern v. Marshall}, 564 U.S. 462 (2011)} (barring non-Article III bankruptcy judges from entering final judgment on common-law counterclaims) and \href{https://supreme.justia.com/cases/federal/us/603/22-859/}{\textit{SEC v. Jarkesy}, 603 U.S. 109 (2024)} (requiring Article III jury adjudication for SEC civil-penalty fraud claims). Through 2026, litigants challenging political review of agency adjudications and non-Article III tribunals continue invoking \textit{Hayburn's Case}, \textit{Plaut}, and \textit{Stern} to police the boundary between executive administration and judicial finality.
\end{itemize}

\section{Pedagogical Critiques and Exam Hypotheticals}

\textbf{Pedagogical Critiques.}
\begin{itemize}
  \item \textbf{The Irony of the Advisory-Opinion Bar:} Legal historians often note a striking paradox: in protesting that Article III courts cannot render advisory opinions, the New York and North Carolina circuit judges wrote unsolicited constitutional letters directly to President Washington outside any pending lawsuit---effectively issuing extrajudicial advisory opinions to condemn advisory opinions.
  \item \textbf{Formalism of the ``Commissioner'' Fiction:} Allowing sitting Article III judges to adjourn court for five minutes, remove their judicial robes, and adjudicate the exact same pension claims in the same courtroom as ``commissioners'' (as Jay and Cushing did, and as \textit{Mistretta} later analogized) invites skepticism about whether institutional independence is truly protected by formal title alone.
  \item \textbf{Doctrinal Expansion Beyond the Record:} Scholars (including Maeva Marcus and Robert Teir) point out that the Supreme Court's only actual ruling in \textit{Hayburn's Case} was a 3--3 tie on Attorney General Randolph's \textit{ex officio} authority; later Courts elevated Dallas's footnote of lower-court letters into black-letter Article III canon.
\end{itemize}

\textbf{Exam Hypothetical 1 (Direct \textit{Hayburn's Case} / \textit{Plaut} Violation).} Congress enacts the ``Cyber-Compensation Act,'' directing federal district courts to adjudicate private claims for federal cyber-attack relief and enter monetary awards, but providing that the Secretary of Homeland Security may ``review and vacate any district court award that the Secretary deems inconsistent with national fiscal priorities.'' \textit{Analysis:} Unconstitutional under \textit{Hayburn's Case} and \textit{Plaut v. Spendthrift Farm}. Congress may assign public-benefit claims to an executive agency in the first instance (with subsequent Article III judicial review), but it cannot vest an Article III court with adjudicatory authority while making the court's judgment subject to executive revision or veto.

\textbf{Exam Hypothetical 2 (``Opposite Side'' Contrast---Prospective Relief and Individual Extrajudicial Service).} By contrast, suppose Congress enacts two separate measures: (a) a statute amending substantive federal environmental law prospectively and directing federal courts to terminate ongoing \textit{prospective} injunctions that no longer satisfy the new statutory standard; and (b) a statute creating an independent advisory commission on federal judicial ethics on which three federal appellate judges may voluntarily serve in their individual capacities without adjudicating cases. \textit{Analysis:} Both provisions are constitutional. Under \href{https://www.law.cornell.edu/supremecourt/text/99-224}{\textit{Miller v. French}, 530 U.S. 327 (2000)}, modifying or terminating \textit{prospective} equitable decrees based on changed underlying law does not retroactively reopen final money judgments in violation of \textit{Hayburn's Case} or \textit{Plaut}. Likewise, under \textit{Mistretta} and the New York circuit judges' ``commissioner'' distinction in \textit{Hayburn's Case}, voluntary extrajudicial service by individual judges on a non-adjudicatory commission does not assign non-judicial duties to Article III courts or compromise judicial finality.

\end{document}
